% ===========================================================================
% VERSIÓN EN ESPAÑOL (solo para lectura y estudio del equipo).
% La entrega oficial es la versión en inglés: main.tex
% Usa los mismos números (generated/) y figuras (figures/) que main.tex.
% ===========================================================================
\documentclass[conference]{IEEEtran}
\IEEEoverridecommandlockouts

\usepackage[utf8]{inputenc}
\usepackage[T1]{fontenc}
\usepackage{cite}
\usepackage{amsmath,amssymb}
\usepackage{graphicx}
\usepackage{xcolor}
\usepackage{url}

\newcommand{\todo}[1]{\textcolor{red}{\textbf{[PENDIENTE: #1]}}}
\newcommand{\figph}[3][\columnwidth]{%
  \IfFileExists{figures/#3.png}{%
    \includegraphics[width=#1]{figures/#3.png}}{%
    \fbox{\parbox[c][2.6cm][c]{\dimexpr#1-2\fboxsep-2\fboxrule\relax}{\centering\small
    \textbf{Insertar figura:} #2\\[2pt] \texttt{figures/#3.png}}}}}
% Generated by workshop3/src/make_report_assets.py -- do not edit by hand
\newcommand{\RandMeanReturn}{23.7}
\newcommand{\RandStdReturn}{12.2}
\newcommand{\RandShapedReturn}{23.7}
\newcommand{\RandMinLen}{9}
\newcommand{\RandMedianLen}{19.0}
\newcommand{\RandMaxLen}{63}
\newcommand{\RandLeftPct}{48.0}
\newcommand{\RandRightPct}{52.0}
\newcommand{\RandThetaFailPct}{100}
\newcommand{\ParityDiff}{0}
\newcommand{\BestConfig}{reference}
\newcommand{\BestSeed}{99}
\newcommand{\BestEpisode}{280}
\newcommand{\BestValReturn}{500.0}
\newcommand{\BestRunSteps}{36{,}500}
\newcommand{\NParams}{17{,}410}
\newcommand{\BestCfgBestVal}{500.0}
\newcommand{\BestCfgFinalVal}{187.9}
\newcommand{\TrainEpisodes}{300}
\newcommand{\MeanRunTime}{69}
\newcommand{\MinFinalVal}{96.2}
\newcommand{\MaxFinalVal}{207.1}
\newcommand{\StepsHalfMean}{6{,}246}
\newcommand{\StepsSolvedMean}{9{,}178}
\newcommand{\RunsSolved}{1/3}
\newcommand{\WOneDatasetSteps}{35{,}874}
\newcommand{\WOneTrainRows}{24{,}782}
\newcommand{\BestCheckpointSteps}{15{,}674}
\newcommand{\TestRandMean}{20.7}
\newcommand{\TestRandStd}{9.2}
\newcommand{\TestRandSuccess}{0}
\newcommand{\TestRandDanger}{24.64}
\newcommand{\TestRandDangerWTwo}{24.59}
\newcommand{\TestRandTheta}{4.01}
\newcommand{\TestRandAbsX}{0.05}
\newcommand{\TestRandXFail}{0}
\newcommand{\NoisyRandMean}{21.0}
\newcommand{\NoisyRandSuccess}{0}
\newcommand{\NoisyRandXFail}{0}
\newcommand{\NoisyRandThetaFail}{100}
\newcommand{\NoisyRandDanger}{23.89}
\newcommand{\NoisyRandDangerWTwo}{23.18}
\newcommand{\TestAltMean}{36.0}
\newcommand{\TestAltStd}{13.5}
\newcommand{\TestAltSuccess}{0}
\newcommand{\TestAltDanger}{14.89}
\newcommand{\TestAltDangerWTwo}{14.53}
\newcommand{\TestAltTheta}{3.95}
\newcommand{\TestAltAbsX}{0.05}
\newcommand{\TestAltXFail}{0}
\newcommand{\NoisyAltMean}{34.1}
\newcommand{\NoisyAltSuccess}{0}
\newcommand{\NoisyAltXFail}{0}
\newcommand{\NoisyAltThetaFail}{100}
\newcommand{\NoisyAltDanger}{15.18}
\newcommand{\NoisyAltDangerWTwo}{15.30}
\newcommand{\TestAngleMean}{42.1}
\newcommand{\TestAngleStd}{9.8}
\newcommand{\TestAngleSuccess}{0}
\newcommand{\TestAngleDanger}{18.93}
\newcommand{\TestAngleDangerWTwo}{10.17}
\newcommand{\TestAngleTheta}{4.74}
\newcommand{\TestAngleAbsX}{0.07}
\newcommand{\TestAngleXFail}{0}
\newcommand{\NoisyAngleMean}{40.8}
\newcommand{\NoisyAngleSuccess}{0}
\newcommand{\NoisyAngleXFail}{0}
\newcommand{\NoisyAngleThetaFail}{100}
\newcommand{\NoisyAngleDanger}{19.21}
\newcommand{\NoisyAngleDangerWTwo}{10.34}
\newcommand{\TestPDMean}{500.0}
\newcommand{\TestPDStd}{0.0}
\newcommand{\TestPDSuccess}{100}
\newcommand{\TestPDDanger}{0.00}
\newcommand{\TestPDDangerWTwo}{0.00}
\newcommand{\TestPDTheta}{0.22}
\newcommand{\TestPDAbsX}{0.37}
\newcommand{\TestPDXFail}{0}
\newcommand{\NoisyPDMean}{464.0}
\newcommand{\NoisyPDSuccess}{64}
\newcommand{\NoisyPDXFail}{36}
\newcommand{\NoisyPDThetaFail}{0}
\newcommand{\NoisyPDDanger}{0.28}
\newcommand{\NoisyPDDangerWTwo}{0.17}
\newcommand{\TestWOneMean}{487.3}
\newcommand{\TestWOneStd}{34.3}
\newcommand{\TestWOneSuccess}{84}
\newcommand{\TestWOneDanger}{0.24}
\newcommand{\TestWOneDangerWTwo}{0.05}
\newcommand{\TestWOneTheta}{3.18}
\newcommand{\TestWOneAbsX}{0.53}
\newcommand{\TestWOneXFail}{16}
\newcommand{\NoisyWOneMean}{466.9}
\newcommand{\NoisyWOneSuccess}{68}
\newcommand{\NoisyWOneXFail}{33}
\newcommand{\NoisyWOneThetaFail}{0}
\newcommand{\NoisyWOneDanger}{0.63}
\newcommand{\NoisyWOneDangerWTwo}{0.08}
\newcommand{\TestWTwoMean}{500.0}
\newcommand{\TestWTwoStd}{0.0}
\newcommand{\TestWTwoSuccess}{100}
\newcommand{\TestWTwoDanger}{0.00}
\newcommand{\TestWTwoDangerWTwo}{0.00}
\newcommand{\TestWTwoTheta}{1.92}
\newcommand{\TestWTwoAbsX}{0.52}
\newcommand{\TestWTwoXFail}{0}
\newcommand{\NoisyWTwoMean}{499.2}
\newcommand{\NoisyWTwoSuccess}{99}
\newcommand{\NoisyWTwoXFail}{1}
\newcommand{\NoisyWTwoThetaFail}{0}
\newcommand{\NoisyWTwoDanger}{0.00}
\newcommand{\NoisyWTwoDangerWTwo}{0.00}
\newcommand{\TestDQNMean}{498.1}
\newcommand{\TestDQNStd}{11.7}
\newcommand{\TestDQNSuccess}{97}
\newcommand{\TestDQNDanger}{10.82}
\newcommand{\TestDQNDangerWTwo}{9.54}
\newcommand{\TestDQNTheta}{6.07}
\newcommand{\TestDQNAbsX}{0.42}
\newcommand{\TestDQNXFail}{3}
\newcommand{\NoisyDQNMean}{370.9}
\newcommand{\NoisyDQNSuccess}{49}
\newcommand{\NoisyDQNXFail}{4}
\newcommand{\NoisyDQNThetaFail}{47}
\newcommand{\NoisyDQNDanger}{11.68}
\newcommand{\NoisyDQNDangerWTwo}{9.44}
\newcommand{\SeedsTestMean}{487.6}
\newcommand{\SeedsTestStd}{16.2}
\newcommand{\SeedsNoisyMean}{450.3}
\newcommand{\SeedsNoisyStd}{56.7}
\newcommand{\SeedsTestSuccess}{95}
\newcommand{\SeedsNoisySuccess}{80}
\newcommand{\NoiseStd}{2}
\newcommand{\WOneThreshold}{0.775}
\newcommand{\WTwoThreshold}{0.892}
\newcommand{\LR}{0.001}
\newcommand{\DiscountGamma}{0.99}
\newcommand{\EpsDecay}{0.99}
\newcommand{\EpsMin}{0.01}
\newcommand{\Buffer}{50{,}000}
\newcommand{\Batch}{64}
\newcommand{\LearnStarts}{1{,}000}
\newcommand{\TargetUpdate}{500}
\newcommand{\GradClip}{10}
\newcommand{\ShapingK}{1}
\newcommand{\EvalEvery}{10}
\newcommand{\WOneFOne}{0.860}
\newcommand{\WOneRecall}{0.886}
\newcommand{\WOnePrecision}{0.834}
\newcommand{\WOnePRAUC}{0.935}
\newcommand{\WTwoFOne}{0.865}
\newcommand{\WTwoRecall}{0.924}
\newcommand{\WTwoPrecision}{0.812}
\newcommand{\WTwoPRAUC}{0.956}
\newcommand{\WTwoWindow}{16}

\renewcommand{\tablename}{TABLA}
\renewcommand{\figurename}{Fig.}
\renewcommand{\abstractname}{Resumen}

\begin{document}

\title{Aprendizaje por Refuerzo para Evitar Fallos en CartPole: un Agente DQN Basado en la Predicción Supervisada de Fallos}

\author{\IEEEauthorblockN{Miguel Angel Torres Ramirez}
\IEEEauthorblockA{Dept. de Ingeniería de Computación\\
Universidad Distrital Francisco Jos\'e de Caldas\\
Email: migatorresr@udistrital.edu.co}
\and
\IEEEauthorblockN{Cristian David Romero Gil}
\IEEEauthorblockA{Dept. de Ingeniería de Computación\\
Universidad Distrital Francisco Jos\'e de Caldas\\
Email: cdromerog@udistrital.edu.co}
\and
\IEEEauthorblockN{Omar Esteban Pe\~na Sanchez}
\IEEEauthorblockA{Dept. de Ingeniería de Computación\\
Universidad Distrital Francisco Jos\'e de Caldas\\
Email: oepenas@udistrital.edu.co}}

\maketitle

\begin{abstract}
En los talleres anteriores predijimos, a partir de trayectorias etiquetadas de CartPole,
si un episodio estaba a punto de fallar. Este informe convierte el mismo sistema en un
problema de control: un agente empuja el carro a la izquierda o a la derecha en cada paso
y debe mantener el péndulo vertical el mayor tiempo posible. Formulamos la tarea como un
Proceso de Decisión de Markov, implementamos un entorno propio de Gymnasium cuya dinámica
coincide exactamente con CartPole-v1, añadimos un \emph{reward shaping} basado en
potenciales construido a partir del hallazgo del Taller~1 (la distancia a los límites de
fallo predice los fallos) y entrenamos una Deep Q-Network implementada desde cero. Se
entrenaron seis configuraciones, cada una cambiando un hiperparámetro, con tres semillas.
El agente seleccionado sobrevive en promedio \TestDQNMean\ de 500 pasos en 100 estados
iniciales reservados (\TestDQNSuccess\% de éxito), frente a \TestRandMean\ de una política
aleatoria, pero no supera a una regla PD diseñada a mano y opera más cerca de los límites
de fallo: los modelos de los Talleres~1 y~2 marcan como peligrosos el \TestDQNDanger\% y el
\TestDQNDangerWTwo\% de sus pasos, y su retorno cae a \NoisyDQNMean\ con ruido en el
actuador. El controlador más robusto fue la LSTM del Taller~2 usada como política con
look-ahead (\NoisyWTwoMean\ pasos con ruido), un modelo supervisado que, a diferencia de
DQN, se apoya en etiquetas generadas por un experto y en un modelo de la dinámica. Las curvas de
entrenamiento muestran un aprendizaje rápido seguido de una degradación de la política en
todas las configuraciones, que diagnosticamos y mitigamos seleccionando el checkpoint con
validación.
\end{abstract}

\begin{IEEEkeywords}
Aprendizaje por refuerzo, Deep Q-Network, proceso de decisión de Markov, reward shaping, CartPole, Gymnasium.
\end{IEEEkeywords}

%% ======================================================================
\section{Introducción}
Los Talleres~1 y~2 trataron los fallos de CartPole como un problema supervisado: dado el
estado del carro y del péndulo, decidir si el episodio fallará en los próximos cinco pasos.
Un clasificador de ese tipo dice \emph{cuándo} se acerca un fallo, pero no \emph{qué hacer}
para evitarlo. Elegir acciones es otro problema: cada empujón cambia los estados que se
pueden alcanzar después, así que la calidad de una decisión solo se conoce cuando sus
consecuencias se desarrollan. Ese es el escenario del aprendizaje por refuerzo (RL), donde
un agente aprende una política interactuando con un entorno y recibiendo recompensas
\cite{sutton2018}.

El RL encaja de forma natural con nuestro dominio por tres razones. Primero, la tarea es una
secuencia de decisiones dependientes. Segundo, el éxito se mide con una recompensa (el
tiempo que el péndulo se mantiene arriba) y no con etiquetas: nuestro dataset contiene las
acciones de una política experta con ruido, no la acción \emph{correcta} para cada estado.
Tercero, el sistema se puede simular, así que las interacciones son baratas.

Las contribuciones de este informe son: (i) una especificación completa del MDP de la
tarea de control, basada en los hallazgos del Taller~1; (ii) un entorno de Gymnasium
validado contra CartPole-v1 y con una política aleatoria; (iii) un agente DQN
\cite{mnih2015} implementado desde cero, con una búsqueda de hiperparámetros de seis
configuraciones y tres semillas; y (iv) una comparación en el conjunto de prueba contra
seis líneas base que no requieren entrenamiento de RL, dos de las cuales convierten en
políticas el clasificador del Taller~1 y la LSTM del Taller~2, junto con un análisis de
eficiencia de muestras.

La Sección~II define el MDP, la Sección~III el entorno, la Sección~IV el agente, la
Sección~V la búsqueda de hiperparámetros, la Sección~VI la evaluación de prueba, la
Sección~VII la eficiencia de muestras, la Sección~VIII discute los resultados y la
Sección~IX concluye.

%% ======================================================================
\section{Formulación del MDP}

\subsection{Planteamiento del Problema}
El agente controla un carro que se mueve sobre un riel de 4.8~m y lleva un péndulo
articulado encima. En cada paso de tiempo (0.02~s) aplica una fuerza horizontal de
$\pm 10$~N. El objetivo es mantener el péndulo vertical y el carro dentro del riel el mayor
tiempo posible, hasta 500 pasos (10~s). La Tabla~\ref{tab:mdp} resume los cinco componentes
del MDP.

\begin{table}[t]
\caption{ESPECIFICACIÓN DEL MDP}
\label{tab:mdp}
\centering
\small
\begin{tabular}{|p{1.5cm}|p{5.9cm}|}
\hline
\textbf{Componente} & \textbf{Definición} \\ \hline
Estado $s_t$ & $(x, \dot x, \theta, \dot\theta) \in \mathbb{R}^4$: posición del carro, velocidad del carro, ángulo del péndulo, velocidad angular \\ \hline
Acciones $a_t$ & $\{0, 1\}$: empujar a la izquierda ($-10$~N) o a la derecha ($+10$~N) \\ \hline
Transición & ecuaciones de movimiento del carro-péndulo, paso de Euler explícito de 0.02~s; determinista dado $(s_t, a_t)$; estado inicial aleatorio $s_0 \sim U(-0.05, 0.05)^4$ \\ \hline
Recompensa & $r_t = 1$ por paso sobrevivido, más un shaping opcional $F = \gamma\Phi(s_{t+1}) - \Phi(s_t)$ \\ \hline
Terminación & fallo si $|\theta| > 12^\circ$ o $|x| > 2.4$~m (terminal); truncamiento a los 500 pasos (no terminal) \\ \hline
\end{tabular}
\end{table}

\subsection{Espacio de Estados}
La observación es el estado físico de 4 dimensiones. Es markoviano: posición, velocidad,
ángulo y velocidad angular determinan por completo la evolución futura del sistema para una
secuencia dada de fuerzas. No le dimos al agente las 16 características diseñadas del
Taller~1: son funciones del estado actual y de los anteriores, así que no aportan
información a un controlador markoviano, y la red Q puede aprender sus propias
características. Las entradas se dividen por escalas físicas fijas (los límites de fallo
para $x$ y $\theta$, y magnitudes típicas para las velocidades); estas constantes vienen de
la definición del entorno, no de los datos.

\subsection{Conjunto de Acciones}
El espacio de acciones es discreto con dos acciones: 0 = empujar a la izquierda y 1 =
empujar a la derecha. No hay acción de ``no hacer nada'', como en el problema original
\cite{barto1983}: el equilibrio resulta de alternar empujones.

\subsection{Dinámica de Transición}
El siguiente estado $s_{t+1} = f(s_t, a_t)$ sigue las ecuaciones del carro-péndulo sin
fricción \cite{barto1983}. Con $F = \pm 10$~N, $m_p = 0.1$~kg, $M = m_c + m_p = 1.1$~kg,
media longitud $l = 0.5$~m y $g = 9.8$~m/s$^2$:
\begin{align}
\ddot\theta &= \frac{g\sin\theta - \cos\theta\,\big(F + m_p l \dot\theta^2 \sin\theta\big)/M}
                    {l\left(\frac{4}{3} - \frac{m_p\cos^2\theta}{M}\right)}, \\
\ddot x &= \frac{F + m_p l\left(\dot\theta^2\sin\theta - \ddot\theta\cos\theta\right)}{M},
\end{align}
seguidas de una actualización de Euler explícita de las cuatro variables con
$\tau = 0.02$~s. La transición es determinista dada la acción. La aleatoriedad viene del
estado inicial y, en la prueba de robustez de la Sección~VI, del ruido gaussiano añadido a
la fuerza.

\subsection{Función de Recompensa}
La \emph{recompensa de la tarea} es $r_t = +1$ por cada paso que el sistema sobrevive, así
que el retorno sin descuento es igual al número de pasos sobrevividos (máximo 500). Es densa
en el tiempo pero no dice nada sobre \emph{qué tan cerca} está el sistema de fallar: la única
señal negativa es la pérdida de recompensas futuras cuando termina el episodio.

El Taller~1 mostró que la distancia a los límites de fallo es el mejor predictor de un
fallo: añadirla subió el ROC-AUC del clasificador lineal de 0.542 a 0.981. Codificamos ese
hallazgo como una función de potencial
\begin{equation}
\Phi(s) = -\frac{k}{2}\left[\left(\frac{x}{2.4}\right)^2 + \left(\frac{\theta}{0.2095}\right)^2\right],
\end{equation}
que vale 0 en la posición vertical y centrada y disminuye hacia los límites, y entrenamos
con la recompensa con shaping $r_t + F_t$, donde $F_t = \gamma\Phi(s_{t+1}) - \Phi(s_t)$ y
$\Phi(\text{terminal}) = 0$. Está demostrado que el shaping basado en potenciales no cambia
la política óptima \cite{ng1999}: a lo largo de un episodio los términos descontados se
cancelan en cadena (suma telescópica) y quedan aproximadamente en $-\Phi(s_0)$, una
constante cercana a cero, así que solo pueden cambiar \emph{qué tan rápido} aprende el
agente, no \emph{qué} aprende. Usamos $k = \ShapingK$; la configuración
\texttt{no\_shaping} ($k = 0$) mide su efecto. Todos los agentes se evalúan solo con la
recompensa de la tarea. El horizonte relevante es largo: con $\gamma = \DiscountGamma$ el
horizonte efectivo $1/(1-\gamma)$ es de unos 100 pasos, comparable al tiempo que tarda un
carro a la deriva en salirse del riel.

\subsection{Terminación del Episodio}
Un episodio \emph{termina} (fallo) cuando $|\theta| > 12^\circ$ o $|x| > 2.4$~m, los mismos
límites que definieron la etiqueta de peligro del Taller~1. Se \emph{trunca} a los 500
pasos: es un límite de tiempo, no un fallo, así que el agente sigue haciendo bootstrap
desde el último estado; solo las terminaciones ponen en cero el valor del siguiente estado.
Llegar al límite de tiempo es un éxito.

%% ======================================================================
\section{Implementación y Validación del Entorno}

\subsection{Implementación}
\texttt{CartPoleControlEnv} (\texttt{workshop3/src/env.py}) hereda de
\texttt{gymnasium.Env} \cite{towers2024}. Declara \texttt{action\_space = Discrete(2)} y un
espacio de observación \texttt{Box} de 4 dimensiones, implementa \texttt{reset()} (estado
inicial aleatorio desde el generador con semilla) y \texttt{step()} con la transición, la
recompensa (de la tarea más shaping) y la lógica de terminación y truncamiento. El
diccionario \texttt{info} reporta la recompensa de la tarea y el desenlace (\emph{fallo por
ángulo}, \emph{fallo por posición} o \emph{límite de tiempo}). \texttt{render()} devuelve una
línea de texto o una imagen RGB dibujada con NumPy (Fig.~\ref{fig:render}).

\begin{figure}[t]
\centering
\figph{cuadros producidos por env.render()}{fig-render}
\caption{Cuadros producidos por \texttt{render()} a lo largo de un episodio de prueba del agente seleccionado (marcas rojas: límites de fallo $|x| = 2.4$~m).}
\label{fig:render}
\end{figure}

\subsection{Validación}
Antes de entrenar se hicieron tres comprobaciones. (1) \texttt{check\_env} de Gymnasium
confirmó la API (espacios, firmas y tipos de datos). (2) Nuestro entorno y CartPole-v1 de
Gymnasium, reiniciados con las mismas semillas y con las mismas acciones aleatorias durante
20 episodios, dieron una diferencia máxima de estado de \ParityDiff: la dinámica es idéntica
a la del sistema que generó nuestro dataset. (3) Ejecutamos 100 episodios con una política
aleatoria uniforme (Fig.~\ref{fig:random}): el retorno medio fue
\RandMeanReturn~$\pm$~\RandStdReturn\ (longitud entre \RandMinLen\ y \RandMaxLen\ pasos,
mediana \RandMedianLen), aparecieron ambas acciones (\RandLeftPct\% izquierda,
\RandRightPct\% derecha) y el \RandThetaFailPct\% de los episodios terminó porque el péndulo
cayó. Con $k = 1$ el retorno medio con shaping fue \RandShapedReturn, casi igual al retorno
de la tarea, como se espera por la propiedad telescópica.

\begin{figure*}[t]
\centering
\figph[0.8\textwidth]{diagnóstico de la política aleatoria}{fig-random-policy}
\caption{Diagnóstico de la política aleatoria en 100 episodios: (a) distribución de la longitud de los episodios, (b) frecuencia de cada acción.}
\label{fig:random}
\end{figure*}

%% ======================================================================
\section{Agente: Deep Q-Network}

\subsection{Algoritmo y Arquitectura}
DQN \cite{mnih2015} es adecuado para un espacio de acciones discreto con dinámica
markoviana. Aproxima la función acción-valor $Q(s, a; w)$, el retorno descontado esperado de
tomar la acción $a$ en el estado $s$ y actuar de forma voraz después, y se entrena hacia el
objetivo de Q-learning \cite{watkins1992}
\begin{equation}
y = r + \gamma\,(1 - d)\max_{a'} Q(s', a'; w^-),
\end{equation}
donde $d = 1$ solo para un fallo y $w^-$ son los pesos de una red objetivo. La red Q
(implementada en PyTorch \cite{paszke2019}) es un perceptrón multicapa
$4 \to 128 \to 128 \to 2$ con activaciones ReLU y salida lineal, con \NParams\ parámetros
entrenables.

\subsection{Ciclo de Entrenamiento}
Cada paso del entorno guarda la transición $(s, a, r, s', d)$ en un buffer de repetición
(replay buffer) de \Buffer\ transiciones; después de \LearnStarts\ pasos, en cada paso se
muestrea uniformemente un mini-lote de \Batch\ transiciones, lo que rompe la correlación
entre estados consecutivos. La pérdida es la de Huber entre $Q(s, a; w)$ e $y$, minimizada
con Adam; la norma del gradiente se recorta en \GradClip\ y la red objetivo es una copia de
la red en línea cada \TargetUpdate\ pasos de gradiente. La exploración es
$\epsilon$-greedy: $\epsilon$ empieza en 1 y se multiplica por el factor de decaimiento al
final de cada episodio, hasta \EpsMin. Cada corrida dura \TrainEpisodes\ episodios; el
retorno, la longitud, $\epsilon$ y la pérdida de cada episodio se registran en TensorBoard y
en archivos CSV en \texttt{runs/workshop3/}. La Tabla~\ref{tab:hparams} lista los valores de
referencia.

\begin{table}[t]
\caption{HIPERPARÁMETROS DE REFERENCIA}
\label{tab:hparams}
\centering
\small
\begin{tabular}{|l|l|}
\hline
\textbf{Hiperparámetro} & \textbf{Valor} \\ \hline
Tasa de aprendizaje (Adam) & \LR \\ \hline
Factor de descuento $\gamma$ & \DiscountGamma \\ \hline
Decaimiento de $\epsilon$ / mínimo & \EpsDecay\ / \EpsMin \\ \hline
Weight decay (L2) & 0 \\ \hline
Intensidad del shaping $k$ & \ShapingK \\ \hline
Replay buffer / tamaño de lote & \Buffer\ / \Batch \\ \hline
Inicio del aprendizaje & \LearnStarts\ pasos \\ \hline
Actualización de la red objetivo & cada \TargetUpdate\ actualizaciones \\ \hline
Recorte del gradiente (norma L2) & \GradClip \\ \hline
Episodios / intervalo de validación & \TrainEpisodes\ / \EvalEvery\ episodios \\ \hline
\end{tabular}
\end{table}

\subsection{Validación y Checkpoints}
Los estados iniciales se dividen en tres conjuntos disjuntos: los episodios de
entrenamiento usan semillas mayores a $10^5$, la validación usa 10 semillas fijas y la
prueba otras 100 semillas fijas. Cada \EvalEvery\ episodios se evalúa al agente de forma
voraz ($\epsilon = 0$) en los episodios de validación. Reportamos el retorno medio de
validación y su media móvil sobre las últimas tres evaluaciones (retorno de validación
suavizado), y guardamos el checkpoint con mayor retorno de validación de cada corrida. La
mejor corrida de la configuración seleccionada se guarda como
\texttt{checkpoints/agent\_best.pt}.

%% ======================================================================
\section{Búsqueda de Hiperparámetros}

\subsection{Protocolo}
Se entrenaron seis configuraciones con las semillas 42, 7 y 99 (18 corridas). Cada una
cambia un solo factor respecto a \texttt{reference}, para poder atribuir las diferencias a
ese factor: una tasa de aprendizaje menor, un factor de descuento menor, un decaimiento de
la exploración más lento, weight decay L2 y sin reward shaping. La configuración se
selecciona solo con resultados de validación: la mayor media (entre semillas) del mejor
retorno de validación; los empates se resuelven con el retorno final de validación (media de
las últimas cinco evaluaciones), que mide la estabilidad al final del entrenamiento.

\begin{table*}[t]
\caption{BÚSQUEDA DE HIPERPARÁMETROS: RESULTADOS DE VALIDACIÓN (MEDIA $\pm$ DESVIACIÓN ESTÁNDAR EN 3 SEMILLAS)}
\label{tab:sweep}
\centering
\small
\setlength{\tabcolsep}{4pt}
\begin{tabular}{|l|c|c|c|c|c|c|c|c|c|}
\hline
\textbf{Configuración} & \textbf{LR} & $\boldsymbol{\gamma}$ & \textbf{Decaim. $\epsilon$} & \textbf{L2} & $\boldsymbol{k}$ &
\textbf{Mejor ret. val.} & \textbf{Ret. val. final} & \textbf{Corridas $\geq$ 475} & \textbf{Pasos a 475} \\ \hline
\textbf{reference} & 0.001 & 0.99 & 0.99 & 0 & 1 & $500.0 \pm 0.0$ & $187.9 \pm 98.3$ & 1/3 & 9{,}178 \\ \hline
lr\_1e-4 & 0.0001 & 0.99 & 0.99 & 0 & 1 & $389.9 \pm 83.3$ & $129.1 \pm 55.4$ & 0/3 & -- \\ \hline
gamma\_0.95 & 0.001 & 0.95 & 0.99 & 0 & 1 & $500.0 \pm 0.0$ & $96.2 \pm 3.0$ & 2/3 & 16{,}902 \\ \hline
eps\_decay\_0.995 & 0.001 & 0.99 & 0.995 & 0 & 1 & $464.4 \pm 50.3$ & $204.2 \pm 72.1$ & 0/3 & -- \\ \hline
weight\_decay\_1e-4 & 0.001 & 0.99 & 0.99 & 0.0001 & 1 & $496.0 \pm 5.7$ & $207.1 \pm 62.4$ & 0/3 & -- \\ \hline
no\_shaping & 0.001 & 0.99 & 0.99 & 0 & 0 & $498.8 \pm 1.6$ & $192.4 \pm 20.2$ & 0/3 & -- \\ \hline

\end{tabular}
\end{table*}

\begin{figure*}[t]
\centering
\figph[0.95\textwidth]{retorno de validación vs pasos del entorno para cada configuración y semilla}{fig-training-curves}
\caption{Retorno de validación (10 episodios, política voraz) en función de los pasos del entorno para cada configuración y semilla. Línea punteada: 475, el umbral para considerar resuelto CartPole-v1.}
\label{fig:curves}
\end{figure*}

\subsection{Resultados}
La Tabla~\ref{tab:sweep} y la Fig.~\ref{fig:curves} resumen las 18 corridas. ``Corridas
$\geq$ 475'' cuenta las corridas cuyo retorno de validación \emph{suavizado} llegó a 475, lo
que exige tres evaluaciones buenas seguidas. Todas las configuraciones muestran un
aprendizaje medible: el retorno de validación sube de unos 10 a varios cientos en los
primeros miles de pasos. La configuración seleccionada es \texttt{\BestConfig}, cuyas tres
semillas alcanzaron un retorno de validación de \BestCfgBestVal; su mejor corrida (semilla
\BestSeed, episodio \BestEpisode) es \texttt{agent\_best.pt}.

El patrón más llamativo lo comparten todas las configuraciones: el retorno final de
validación (\MinFinalVal\ a \MaxFinalVal) está muy por debajo del mejor. Los agentes aprenden
una buena política y luego la pierden, una inestabilidad conocida de DQN causada por hacer
bootstrap desde un objetivo que se mueve, con aproximación de funciones. La
Fig.~\ref{fig:explore}(b) muestra cómo ocurre en la configuración de referencia: durante los
primeros dos mil pasos los episodios terminan porque cae el péndulo; una vez que el agente
aprende a sostenerlo, casi todos los episodios de validación que no llegan a 500 pasos
terminan porque el carro se sale del riel, incluidos los de las políticas degradadas al
final del entrenamiento (pasados unos 22 mil pasos la curva proviene de una sola semilla, la
única corrida que duró tanto). La política
degradada mantiene el péndulo arriba pero deja que el carro se desplace, un fallo cuyo costo
aparece unos 100 pasos después, en el borde del horizonte efectivo de $\gamma = 0.99$.
Seleccionar los checkpoints con validación es lo que hace que el agente final sea útil pese
a esta degradación.

En cuanto a los factores individuales: con la tasa de aprendizaje menor ninguna semilla
llegó a un retorno de validación de 500, así que su mejor retorno es el más bajo de la
búsqueda; $\gamma = 0.95$ llegó a 500 en todas las semillas pero fue la que más se degradó,
hasta un retorno final de unos 96, coherente con un horizonte de unos 20 pasos que no
alcanza a ver el fallo por deriva; el decaimiento más lento de la exploración y el weight
decay no cambiaron el panorama; y quitar el shaping produjo prácticamente los mismos
retornos mejor y final que la referencia. Con $k = 1$ el shaping no tuvo un efecto medible,
lo cual es compatible con su papel teórico (no puede cambiar la política óptima) pero
significa que tampoco aceleró el aprendizaje.

\begin{figure*}[t]
\centering
\figph[0.95\textwidth]{calendarios de epsilon y cómo terminan los episodios de validación}{fig-exploration}
\caption{(a) Tasa de exploración por episodio para los dos factores de decaimiento. (b) Configuración de referencia: proporción de episodios de validación que terminan por posición del carro o por ángulo del péndulo durante el entrenamiento (media de tres semillas).}
\label{fig:explore}
\end{figure*}

%% ======================================================================
\section{Evaluación de Prueba y Comparación}

\subsection{Protocolo y Líneas Base}
El conjunto de prueba contiene 100 estados iniciales que nunca se usaron para entrenar ni
para seleccionar. Se evalúa en dos condiciones: \emph{nominal} (la dinámica del
entrenamiento) y \emph{con ruido}, con ruido gaussiano en el actuador de desviación estándar
\NoiseStd~N sumado a cada empujón, una perturbación que el agente nunca vio. El agente DQN se
compara con seis políticas que no necesitan entrenamiento de RL:
\begin{itemize}
\item \textbf{Aleatoria (Random):} acciones aleatorias uniformes.
\item \textbf{Alternada (fija):} izquierda, derecha, izquierda\dots, una estrategia fija en lazo abierto.
\item \textbf{Regla de ángulo (voraz):} empujar hacia el lado al que se inclina el péndulo.
\item \textbf{Regla PD (heurística):} empujar hacia el signo de $\theta + 0.5\,\dot\theta$; la ganancia se eligió con los episodios de validación.
\item \textbf{Clasificador W1 (look-ahead):} el clasificador Gradient Boosting del Taller~1 usado como política. Para cada acción, la dinámica conocida da el siguiente estado y el clasificador estima su probabilidad de peligro; se elige la acción con la menor probabilidad. A diferencia de DQN, usa un modelo de la dinámica.
\item \textbf{LSTM W2 (look-ahead):} la misma regla con el modelo del Taller~2, una LSTM de 2 capas que lee los últimos \WTwoWindow\ pasos del episodio (estados y las acciones realmente tomadas). El paso candidato lleva la acción que conduce a él, porque todavía no se sabe qué acción se tomará allí. El modelo se carga desde \texttt{checkpoints/model\_best.pt} sin reentrenarlo; en los episodios de prueba del Taller~2 nuestro puente reproduce exactamente su F1 de prueba reportado.
\end{itemize}

Además del retorno, reportamos la tasa de éxito (llegar a 500 pasos), la proporción de
fallos por posición del carro, el $|\theta|$ medio y dos \emph{tasas de peligro}: la
proporción de pasos que el clasificador del Taller~1 (umbral \WOneThreshold) y la LSTM del
Taller~2 (umbral \WTwoThreshold), cada uno con el umbral elegido en su propio conjunto de
validación, marcan como dentro de los cinco pasos previos a un fallo. Estas métricas conectan
los tres talleres: miden cuánto tiempo pasa cada política en estados que los modelos
supervisados consideran peligrosos.

\begin{table*}[t]
\caption{RESULTADOS DE PRUEBA EN 100 EPISODIOS RESERVADOS (RETORNO = PASOS SOBREVIVIDOS, MÁXIMO 500)}
\label{tab:test}
\centering
\small
\setlength{\tabcolsep}{3.5pt}
\begin{tabular}{|l|c|c|c|c|c|c||c|c|}
\hline
& \multicolumn{6}{c||}{\textbf{Prueba nominal}} & \multicolumn{2}{c|}{\textbf{Con ruido (\NoiseStd~N)}} \\ \hline
\textbf{Política} & \textbf{Retorno} & \textbf{Éxito (\%)} & \textbf{Fallo $x$ (\%)} & \textbf{Pel. W1 (\%)} & \textbf{Pel. W2 (\%)} & \textbf{$|\theta|$ (grados)} & \textbf{Retorno} & \textbf{Éxito (\%)} \\ \hline
Random & $20.7 \pm 9.2$ & 0 & 0 & 24.64 & 24.59 & 4.01 & $21.0 \pm 11.5$ & 0 \\ \hline
Alternating (fixed) & $36.0 \pm 13.5$ & 0 & 0 & 14.89 & 14.53 & 3.95 & $34.1 \pm 14.2$ & 0 \\ \hline
Angle rule (greedy) & $42.1 \pm 9.8$ & 0 & 0 & 18.93 & 10.17 & 4.74 & $40.8 \pm 10.7$ & 0 \\ \hline
PD rule (heuristic) & $500.0 \pm 0.0$ & 100 & 0 & 0.00 & 0.00 & 0.22 & $464.0 \pm 65.0$ & 64 \\ \hline
W1 classifier (look-ahead) & $487.3 \pm 34.3$ & 84 & 16 & 0.24 & 0.05 & 3.18 & $466.9 \pm 56.0$ & 68 \\ \hline
W2 LSTM (look-ahead) & $500.0 \pm 0.0$ & 100 & 0 & 0.00 & 0.00 & 1.92 & $499.2 \pm 8.2$ & 99 \\ \hline
\textbf{DQN (agent\_best.pt)} & $498.1 \pm 11.7$ & 97 & 3 & 10.82 & 9.54 & 6.07 & $370.9 \pm 158.7$ & 49 \\ \hline
DQN, 3 seeds of reference & $487.6 \pm 16.2$ & 95 & -- & -- & -- & -- & $450.3 \pm 56.7$ & 80 \\ \hline

\end{tabular}
\end{table*}

\begin{figure}[t]
\centering
\figph{retorno de prueba de cada política}{fig-test-comparison}
\caption{Retorno medio de prueba ($\pm$ una desviación estándar) de cada política en las condiciones nominal y con ruido.}
\label{fig:test}
\end{figure}

\subsection{Resultados}
La Tabla~\ref{tab:test} y la Fig.~\ref{fig:test} muestran los resultados. En la condición
nominal el agente seleccionado sobrevive \TestDQNMean~$\pm$~\TestDQNStd\ pasos con
\TestDQNSuccess\% de éxito, muy por encima de las políticas aleatoria (\TestRandMean),
alternada (\TestAltMean) y de la regla de ángulo (\TestAngleMean). Entre las tres semillas de
la configuración seleccionada la media es \SeedsTestMean~$\pm$~\SeedsTestStd, así que el
resultado no depende de una semilla afortunada. El agente no superó a la regla PD ni a la
política guiada por W2, que resuelven todos los episodios de prueba nominales, y queda
ligeramente por encima del clasificador W1 usado como política (\TestWOneMean,
\TestWOneSuccess\% de éxito).

La tasa de peligro y las trayectorias revelan \emph{cómo} sobrevive el agente. La regla PD
mantiene el péndulo casi vertical ($|\theta|$ medio de \TestPDTheta$^\circ$) y pasa el
\TestPDDanger\% de sus pasos en peligro. El agente seleccionado hace oscilar el péndulo entre
unos $-11^\circ$ y $+11^\circ$, cerca del límite de $12^\circ$ (Fig.~\ref{fig:traj}). Los dos
modelos supervisados, entrenados por separado y con entradas distintas, coinciden: el
clasificador del Taller~1 marca como peligrosos el \TestDQNDanger\% de sus pasos y la LSTM
del Taller~2 el \TestDQNDangerWTwo\%. Ambas políticas maximizan igual la recompensa de la
tarea, pero el agente tiene un margen de seguridad mucho menor, y la prueba con ruido lo deja
ver: su retorno cae a \NoisyDQNMean\ con \NoisyDQNSuccess\% de éxito, y el
\NoisyDQNThetaFail\% de los episodios termina porque cae el péndulo. La regla PD
(\NoisyPDMean) y la política guiada por W1 (\NoisyWOneMean) se degradan menos, sobre todo
por deriva del carro. La política guiada por W2 es la más robusta de todas:
\NoisyWTwoMean\ pasos con \NoisyWTwoSuccess\% de éxito con ruido, manteniendo el péndulo en
un $|\theta|$ medio de \TestWTwoTheta$^\circ$ en la prueba nominal.

La Tabla~\ref{tab:seeds} muestra que esta fragilidad es del checkpoint seleccionado, no de
la configuración: las otras dos semillas tienen ambas tasas de peligro por debajo de 0.1\%
y conservan la mayor parte de su desempeño con ruido. Entre estos tres agentes, las tasas de
peligro, medidas en la condición nominal, identificaron al frágil, mientras que el retorno de
validación, igual a 500 para los tres, no pudo hacerlo. Lo presentamos como una observación
sobre tres agentes, no como un resultado general.

\begin{table}[t]
\caption{CONFIGURACIÓN SELECCIONADA: RESULTADOS DE PRUEBA POR SEMILLA}
\label{tab:seeds}
\centering
\small
\setlength{\tabcolsep}{3pt}
\begin{tabular}{|l|c|c|c|c|c||c|c|}
\hline
& \multicolumn{5}{c||}{\textbf{Nominal}} & \multicolumn{2}{c|}{\textbf{Con ruido}} \\ \hline
\textbf{Semilla} & \textbf{Ret.} & \textbf{Éxito} & \textbf{P. W1} & \textbf{P. W2} & $|\theta|$ & \textbf{Ret.} & \textbf{Éxito} \\ \hline
42 & 464.7 & 88 & 0.09 & 0.00 & 1.06 & 479.9 & 92 \\ \hline
7 & 500.0 & 100 & 0.08 & 0.00 & 5.26 & 500.0 & 100 \\ \hline
99 (agent\_best) & 498.1 & 97 & 10.82 & 9.54 & 6.07 & 370.9 & 49 \\ \hline

\end{tabular}
\end{table}

\subsection{Comparación con los Talleres 1 y 2}
Los modelos supervisados de los talleres anteriores miden otra cantidad (clasificación de
estados de peligro), así que sus métricas no se pueden poner en la misma columna que un
retorno; la Tabla~\ref{tab:w12} relaciona ambas miradas. En los episodios de prueba comunes
de los Talleres~1 y~2, la LSTM del Taller~2 fue mejor predictor (F1 \WTwoFOne\ frente a
\WOneFOne, recall \WTwoRecall\ frente a \WOneRecall, PR-AUC \WTwoPRAUC\ frente a
\WOnePRAUC), y como política con look-ahead también fue mejor controlador (\TestWTwoMean\
frente a \TestWOneMean\ en nominal, \NoisyWTwoMean\ frente a \NoisyWOneMean\ con ruido).
Una razón plausible, que no probamos por separado, es que su ventana de 16 pasos le permite
juzgar hacia dónde va el sistema, mientras que las características del W1 solo miran tres
pasos atrás. Ambos controladores supervisados se apoyan en dos recursos que el agente DQN no
tuvo: etiquetas producidas por una política experta y un modelo de la dinámica para simular
cada acción.

\begin{table}[t]
\caption{MODELOS SUPERVISADOS COMO PREDICTORES (PRUEBA W1/W2) Y COMO POLÍTICAS CON LOOK-AHEAD (PRUEBA W3)}
\label{tab:w12}
\centering
\small
\setlength{\tabcolsep}{2.5pt}
\begin{tabular}{|l|c|c|c||c|c|}
\hline
& \multicolumn{3}{c||}{\textbf{Predicción}} & \multicolumn{2}{c|}{\textbf{Control (retorno)}} \\ \hline
\textbf{Modelo} & \textbf{F1} & \textbf{Recall} & \textbf{PR-AUC} & \textbf{Nominal} & \textbf{Ruido} \\ \hline
W1 Gradient Boosting & \WOneFOne & \WOneRecall & \WOnePRAUC & \TestWOneMean & \NoisyWOneMean \\ \hline
W2 LSTM & \WTwoFOne & \WTwoRecall & \WTwoPRAUC & \TestWTwoMean & \NoisyWTwoMean \\ \hline
DQN (este trabajo) & -- & -- & -- & \TestDQNMean & \NoisyDQNMean \\ \hline
\end{tabular}
\end{table}

\subsection{Trayectorias de Ejemplo}
La Fig.~\ref{fig:traj} muestra cuatro episodios de prueba del agente seleccionado: ángulo del
péndulo, posición del carro, acciones (empujar a la derecha en azul) y retorno acumulado. El
agente alterna empujones en ráfagas, lo que produce una oscilación regular del péndulo
mientras el carro se mantiene a menos de un metro del centro; el retorno crece en uno por
paso hasta el límite de tiempo.

\begin{figure*}[t]
\centering
\figph[0.95\textwidth]{cuatro episodios de prueba del agente seleccionado}{fig-trajectories}
\caption{Cuatro episodios de prueba del agente seleccionado: ángulo del péndulo, posición del carro, acción tomada en cada paso y retorno acumulado de la tarea.}
\label{fig:traj}
\end{figure*}

%% ======================================================================
\section{Eficiencia de Muestras}
La Fig.~\ref{fig:sample} grafica el retorno de validación suavizado de la configuración
seleccionada contra el número de pasos del entorno. En promedio, sus semillas alcanzaron por
primera vez un retorno suavizado de 195 tras \StepsHalfMean\ pasos; el nivel de resuelto,
475, lo alcanzaron \RunsSolved\ semillas (tras \StepsSolvedMean\ pasos), y los mejores
checkpoints se guardaron en promedio tras \BestCheckpointSteps\ pasos.

\begin{figure*}[t]
\centering
\figph[0.95\textwidth]{curva de aprendizaje de la configuración seleccionada frente a las líneas base}{fig-sample-efficiency}
\caption{Eficiencia de muestras: retorno de validación suavizado de la configuración seleccionada (media y mínimo--máximo de tres semillas) frente a los pasos del entorno, con el retorno de prueba de las líneas base (la regla PD y la LSTM W2 comparten la línea en 500) y el número de pasos del entorno almacenados en el dataset de los Talleres~1 y~2.}
\label{fig:sample}
\end{figure*}

Como comparación, el dataset de los Talleres~1 y~2 contiene \WOneDatasetSteps\ pasos del
entorno, de los cuales se usaron \WOneTrainRows\ filas etiquetadas para entrenar ambos
modelos supervisados. Por tanto, el agente de RL alcanzó un control competente con menos
interacciones con el entorno de las que hicieron falta para construir el dataset supervisado,
aunque los modelos supervisados, convertidos en controladores, terminaron siendo más
robustos. La comparación tiene dos matices. Los modelos
supervisados aprenden a \emph{predecir} fallos, y convertirlos en un controlador requirió el
modelo de la dinámica que usa la política con look-ahead; y el propio dataset lo produjo una
política experta, una fuente de conocimiento que el agente de RL no tuvo. Por otro lado, la
regla PD no necesita ningún dato, solo conocimiento de la física.

%% ======================================================================
\section{Discusión}

\emph{¿El RL superó al aprendizaje supervisado?} No. En el entorno nominal el agente de RL
iguala a la política guiada por W1 y se acerca a la regla PD y a la política guiada por W2,
aprendiendo solo de recompensas y sin un modelo de la dinámica. Con ruido en el actuador fue
el menos robusto de esos cuatro, y la LSTM W2 con look-ahead fue el mejor controlador. La
comparación no es en igualdad de condiciones: los controladores supervisados usaron
etiquetas generadas por un experto y el simulador como modelo de la dinámica, y el agente
DQN solo sus propias interacciones. En un sistema tan simple y tan bien entendido, una regla
diseñada a mano también sigue siendo una línea base fuerte.

\emph{Principales retos.} El principal fue la inestabilidad del entrenamiento: todas las
configuraciones alcanzaron una buena política y luego se degradaron, así que el agente final
depende de seleccionar el checkpoint con validación (Fig.~\ref{fig:curves}). El segundo es que
la recompensa de la tarea no premia los márgenes de seguridad: sobrevivir con el péndulo
cerca del límite gana la misma recompensa que mantenerlo vertical, lo que explica el
comportamiento oscilante del agente seleccionado y su fragilidad con ruido. Validar solo con
episodios nominales no pudo detectarlo; añadir ruido a los episodios de validación, o las
tasas de peligro W1/W2 al criterio de selección, es una mejora natural.

\emph{Exploración frente a explotación.} Con un decaimiento de \EpsDecay\ por episodio,
$\epsilon$ baja de 0.1 tras unos 230 episodios; el decaimiento más lento de 0.995 lo mantiene
por encima de 0.2 durante toda la corrida (Fig.~\ref{fig:explore}a). El decaimiento más lento
produjo episodios de entrenamiento más cortos y ruidosos, y por tanto menos pasos del entorno
por corrida, sin mejorar el mejor retorno de validación. Al final del entrenamiento, con poca
exploración, la política derivó hacia el modo de fallo por posición del carro, lo que sugiere
que la baja exploración junto con el bootstrap deja sin corregir errores en estados poco
visitados (cerca de los extremos del riel).

\emph{Reward shaping.} El shaping basado en potenciales con $k = 1$ no cambió los resultados
de forma medible. Como no puede cambiar la política óptima, tampoco puede hacer que el agente
prefiera márgenes de seguridad mayores; una recompensa que penalice explícitamente la
cercanía a los límites cambiaría la tarea, y por tanto la política óptima, lo cual es una
decisión de diseño y no de ajuste.

\emph{Limitaciones.} Tres semillas por configuración y diez episodios de validación dan
estimaciones ruidosas; los hiperparámetros se variaron de uno en uno, así que no se
exploraron sus interacciones; y la comparación con los modelos supervisados es indirecta
porque resuelven una tarea de predicción. Double DQN, actualizaciones suaves de la red
objetivo o replay con prioridades podrían reducir la degradación.

%% ======================================================================
\section{Conclusiones}
Formulamos el control de CartPole como un MDP basado en nuestro trabajo previo de predicción
de fallos, implementamos y validamos un entorno de Gymnasium con la misma dinámica que
CartPole-v1, y entrenamos un agente DQN hecho desde cero con una búsqueda de hiperparámetros
de seis configuraciones y tres semillas. El agente aprendió de forma medible, sobreviviendo
\TestDQNMean\ de 500 pasos en episodios reservados frente a \TestRandMean\ de una política
aleatoria, pero no superó a una regla PD y opera con un margen de seguridad pequeño que lo
vuelve frágil ante el ruido del actuador. Los modelos de los Talleres~1 y~2, reutilizados
como monitores de peligro, revelaron ese margen, y la LSTM del Taller~2, reutilizada como
política con look-ahead, fue el controlador más robusto. Como trabajo futuro: validación con ruido, funciones
de recompensa que valoren explícitamente los márgenes de seguridad y variantes estabilizadas
de DQN.

\section*{Reproducibilidad}
El código, las configuraciones, los registros y los checkpoints están en
\url{https://github.com/Tonyaa11/Cartpole} (carpeta \texttt{workshop3/}). Todos los números y
figuras de este informe los genera \texttt{workshop3/src/make\_report\_assets.py} a partir de
los registros; las semillas son fijas (entrenamiento 42, 7, 99; validación 10000--10009;
prueba 20000--20099). El entrenamiento de redes neuronales puede variar ligeramente entre
versiones de librerías y hardware.

\begin{thebibliography}{00}
\bibitem{sutton2018} R. S. Sutton and A. G. Barto, \emph{Reinforcement Learning: An Introduction}, 2nd ed. Cambridge, MA, USA: MIT Press, 2018.
\bibitem{mnih2015} V. Mnih et al., ``Human-level control through deep reinforcement learning,'' \emph{Nature}, vol. 518, no. 7540, pp. 529--533, 2015.
\bibitem{watkins1992} C. J. C. H. Watkins and P. Dayan, ``Q-learning,'' \emph{Mach. Learn.}, vol. 8, no. 3--4, pp. 279--292, 1992.
\bibitem{barto1983} A. G. Barto, R. S. Sutton, and C. W. Anderson, ``Neuronlike adaptive elements that can solve difficult learning control problems,'' \emph{IEEE Trans. Syst., Man, Cybern.}, vol. SMC-13, no. 5, pp. 834--846, 1983.
\bibitem{ng1999} A. Y. Ng, D. Harada, and S. Russell, ``Policy invariance under reward transformations: Theory and application to reward shaping,'' in \emph{Proc. 16th Int. Conf. Mach. Learn. (ICML)}, 1999, pp. 278--287.
\bibitem{towers2024} M. Towers et al., ``Gymnasium: A standard interface for reinforcement learning environments,'' arXiv:2407.17032, 2024.
\bibitem{paszke2019} A. Paszke et al., ``PyTorch: An imperative style, high-performance deep learning library,'' in \emph{Adv. Neural Inf. Process. Syst. (NeurIPS)}, vol. 32, 2019.
\end{thebibliography}

\end{document}
